\documentclass[10pt,a4paper]{amsart}
\usepackage[margin=19mm]{geometry}
\usepackage{amsmath,amssymb,mathtools}
\usepackage[scr=rsfs,cal=euler]{mathalfa}
\usepackage{xfrac}
\usepackage[unicode,hidelinks,bookmarks=false]{hyperref}
\newcommand{\ie}{i.e.\ }
\newcommand{\cf}{cf.\ }
\newcommand{\resp}{respectively\ }
\newcommand{\Subsection}{\subsection}
\providecommand{\nobreakdash}{\nobreak\hbox{-}}
\setlength{\emergencystretch}{2em}
\multlinegap0pt
\usepackage{amssymb}
\usepackage{mathtools}
\usepackage{bm}
\usepackage{xcolor}
\newtheorem{lemma}{Lemma}
\newtheorem{proposition}{Proposition}
\newcommand{\C}{\mathbb{C}}
\newcommand{\leaf}[1]{\mathcal{L}_{#1}}
\title{Holomorphic Linear $\C^k$-Actions, Trace Foliations, and Higher-Rank Poincaré Dynamics}
\author{Aubin Arroyo, Carlos Cabrera, Jos\'e Seade,, Alberto Verjovsky}
\date{Source: arXiv:2608.01286v1 (CC BY 4.0)}
\begin{document}
\maketitle

\noindent\emph{Source and licence.} Holomorphic Linear $\C^k$-Actions, Trace Foliations, and Higher-Rank Poincaré Dynamics. Version arXiv:2608.01286v1 (CC BY 4.0), distributed by arXiv under the Creative Commons Attribution 4.0 International licence, http://creativecommons.org/licenses/by/4.0/, as recorded in the arXiv licence metadata for this exact version and shown on the abstract page https://arxiv.org/abs/2608.01286v1. Full licence notice: \texttt{licenses/arxiv-2608.01286-CC-BY-4.0.txt}.\par\smallskip

\noindent\textit{Editorial scope.} Counted unit: the article's proposition foliation (source0.tex lines 1244--1247). The article's complete original proof of it is delivered with it (source0.tex lines 1249--1414, one \texttt{proof} environment of 166 lines). No mathematics is written, repaired or re-derived: the statement and the proof are the article's own lines, in the article's own notation, and no retained line is substituted or re-flowed.  Retained uncounted as the source-local context the proof needs: lines 1180--1189.  Nothing was omitted from any retained span, so the package carries no omission record.  No retained line contains a citation, so the wrapper carries no bibliography and no bibliography entry.  No retained line contains a figure environment, an image-inclusion command or drawing code, so no image is delivered and the package declares no asset dependency.  Editorial formatting only, in addition: an AMS \texttt{amsart} preamble that restates the article's own declarations and macros used by the retained lines, namely the environments \texttt{lemma}, \texttt{proposition} on separate counters, together with the standard packages those lines need.  The retained environments print in this document as Lemma 1, Proposition 1 in order of appearance, whereas the article prints the same items with the numbering of its own full document.  The complete native source of the article as distributed in the arXiv e-print for this exact version is bundled unchanged as \texttt{sources/206613/source0.tex}.
\begin{lemma}[Local product chart]\label{lem:Local product chart} 
	Let $\Sigma_p$ be a Stefan--Sussmann transversal at $p$, and let $q=q(p)$.
	Then, after shrinking $\Sigma_p$ if necessary, there exist a neighborhood $B$ of $0$ in $\C^q$, a neighborhood $V$ of $p$ in $\C^n$, and a biholomorphism 
	\[
		\Psi:B\times\Sigma_p\longrightarrow V
	\]
	such that 
	$\Psi(0,y)=y$ for every $y\in\Sigma_p$.
	Moreover, for every $y\in\Sigma_p$, the slice $\Psi(B\times\{y\})$ is contained in the ambient orbit $\leaf{y}$.
\end{lemma}

\begin{proposition}[Uniqueness of the transverse foliation]\label{Uniqueness transverse foliation}
	Let $\Sigma_1$ and $\Sigma_2$ be two Stefan--Sussmann transversals at a point $p$.
	Then, after shrinking them if necessary, the transverse orbit foliations induced on $\Sigma_1$ and $\Sigma_2$ are biholomorphically equivalent as germs at $p$.
\end{proposition}

\begin{proof}
%	Let $\Psi:B\times\Sigma_1\longrightarrow V$ be a local product chart provided by Lemma~\ref{lem:Local product chart}. Using the local product chart $\Psi$, we identify $\Sigma_1$ with the zero section $\{0\}\times\Sigma_1$.
	Let $\Psi:B\times\Sigma_1\longrightarrow V$ be the local product chart provided by Lemma~\ref{lem:Local product chart}, and identify $\Sigma_1$ with the zero section $\{0\}\times\Sigma_1$.
	Let
	$$
		\widetilde{\Sigma}_2:=\Psi^{-1}(\Sigma_2)\subset B\times\Sigma_1.
	$$

	Since $d\Psi_{(0,p)}$ identifies the tangent space to the fibers of 
	\[
		\pi_2:B\times\Sigma_1\longrightarrow\Sigma_1
	\]
	with $T_p\leaf{p} $, the transversality of $\Sigma_2$ to $\leaf{p}$ implies that $\widetilde{\Sigma}_2$ is transverse at $(0,p)$ to the fibers of $\pi_2$.
	Equivalently, the restriction 
	\[
		\pi_2|_{\widetilde{\Sigma}_2}: \widetilde{\Sigma}_2 \longrightarrow \Sigma_1
	\]
	has invertible differential at $(0,p)$.

%	By the holomorphic inverse function theorem, after shrinking if necessary, $\pi_2|_{\widetilde{\Sigma}_2}$ is a biholomorphism onto a neighborhood of $p$ in $\Sigma_1$.
	After shrinking $\Sigma_1$ and $\widetilde{\Sigma}_2$ if necessary, the holomorphic inverse function theorem implies that $\pi_2|_{\widetilde{\Sigma}_2}$ is a biholomorphism onto a neighborhood of $p$ in $\Sigma_1$.
%	Therefore $\widetilde{\Sigma}_2$ is the graph of a unique holomorphic map 
%	$$ 
%		a:\Sigma_1\longrightarrow B, \qquad a(p)=0.
%	$$
%	Thus
%	$$
%		\widetilde{\Sigma}_2 = \{(a(y),y):y\in\Sigma_1\}.
%	$$
	Therefore, $\widetilde{\Sigma}_2$ is the graph of a unique holomorphic map $a:\Sigma_1\longrightarrow B$, with $a(p)=0$; namely,
	\[
		\widetilde{\Sigma}_2 = \{(a(y),y):y\in\Sigma_1\}.
	\]
	Define the vertical holomorphic vector field 
	\[
		Y = \sum_{j=1}^q a_j(y)\frac{\partial}{\partial t_j}
	\]
	on $B\times\Sigma_1$.
	Observe that $Y(0,p)=0$, so its flow fixes the point $(0,p)$.

	Since $Y$ has no component in the $\Sigma_1$-direction, the coordinate $y$ remains constant along every integral curve of $Y$.
	Consequently, the coefficients $a_j(y)$ remain constant along each integral curve.
%	Hence, if $(t_0,y)$ is the initial point of an integral curve, its solution is 
%	$$ 
%		t(s)=t_0+s\,a(y), \qquad y(s)=y.
%	$$
%	Hence, the integral curve through $(t_0,y)$ is given by $t(s)=t_0+s\,a(y)$ and $y(s)=y$.
%	Therefore the time-one map of the flow of $Y$ is
%	$$
%		(t,y)\longmapsto(t+a(y),y).
%	$$
%	In particular, it sends the zero section
%	$$
%		\{0\}\times\Sigma_1
%	$$
%	biholomorphically onto the graph of $a$, namely onto $\widetilde{\Sigma}_2$.
	Therefore, the time-one map of the flow of $Y$ is given by $(t,y)\mapsto(t+a(y),y)$, and its restriction to the zero section $\{0\}\times\Sigma_1$ is a biholomorphism onto the graph of $a$, which is precisely $\widetilde{\Sigma}_2$.

%	It remains to show that the vector field $\Psi_*Y$ is tangent to the ambient Stefan--Sussmann foliation.
%	Let 
%	\[
%		\Phi_j^s=\exp(sX_j)
%	\]
% 	denote the local flow of $X_j$.
%	Since the action is abelian, the vector fields $X_1,\ldots,X_q$ commute, and therefore so do their flows.

It remains to prove that $\Psi_*Y$ is tangent to the ambient
Stefan--Sussmann foliation. For $j=1,\ldots,q$, let
$\Phi_j^s=\exp(sX_j)$ denote the local flow of $X_j$. Since the action
is abelian, the fundamental vector fields $X_1,\ldots,X_q$ commute,
and hence so do their local flows wherever the corresponding
compositions are defined.

%	Fix $j$.
%	Using the commutativity of the flows, write 
%	$$ 
%		\Psi(t,y) = \Phi_j^{t_j} \circ \Phi_1^{t_1} \circ\cdots\circ \widehat{\Phi_j^{t_j}} \circ\cdots\circ \Phi_q^{t_q}(y).
%	$$
%	If
%	$$
%		z = \Phi_1^{t_1} \circ\cdots\circ \widehat{\Phi_j^{t_j}} \circ\cdots\circ \Phi_q^{t_q}(y),
%	$$
%	then
%	$$
%		\Psi(t,y)=\Phi_j^{t_j}(z).
%	$$
%	Differentiating with respect to $t_j$, we obtain
%	$$
%		\frac{\partial\Psi}{\partial t_j}(t,y) = \frac{d}{ds}\bigg|_{s=t_j}\Phi_j^s(z) = X_j(\Phi_j^{t_j}(z)) = X_j(\Psi(t,y)).
%	$$
%	Equivalently,
%	$$
%		\Psi_*\!\left(\frac{\partial}{\partial t_j}\right)=X_j.
%	$$
 
Fix $j\in\{1,\ldots,q\}$. By commutativity of the flows, we may write
\[
    \Psi(t,y)=\Phi_j^{t_j}(z_j),
    \qquad
    z_j:=
    \Phi_1^{t_1}\circ\cdots\circ
    \widehat{\Phi_j^{t_j}}\circ\cdots\circ
    \Phi_q^{t_q}(y),
\]
where $z_j$ is independent of $t_j$. Therefore,
\[
    \frac{\partial\Psi}{\partial t_j}(t,y)
    =
    \left.\frac{d}{ds}\right|_{s=t_j}\Phi_j^s(z_j)
    =
    X_j\bigl(\Phi_j^{t_j}(z_j)\bigr)
    =
    X_j\bigl(\Psi(t,y)\bigr).
\]
Equivalently,
\[
    \Psi_*\!\left(\frac{\partial}{\partial t_j}\right)=X_j.
\]

 
%	To simplify notation, we continue to denote the coefficient functions $a_j\circ\pi_2\circ\Psi^{-1}$ simply by $a_j$.
%	Thus 
%	\[
%		\Psi_*Y = \sum_{j=1}^q a_jX_j,
%	\]
% 	which is a holomorphic linear combination of fundamental vector fields.
%	Therefore $\Psi_*Y$ is tangent to the Stefan--Sussmann foliation.

To simplify notation, we continue to write $a_j$ for the coefficient
function $a_j\circ\pi_2\circ\Psi^{-1}$. It follows that
\[
    \Psi_*Y=\sum_{j=1}^q a_jX_j.
\]
Since the coefficients $a_j$ are holomorphic and the $X_j$ are
fundamental vector fields, $\Psi_*Y$ is tangent to the ambient
Stefan--Sussmann foliation.

%	Since every integral curve of $\Psi_*Y$ is contained in an ambient leaf, its time-one flow preserves every ambient leaf.
%	Consequently, its restriction defines a biholomorphism of germs 
%	$$ 
%		F:(\Sigma_1,p)\longrightarrow(\Sigma_2,p).
%	$$
%	Moreover, for every ambient leaf  $\leaf{}$ meeting the neighborhood, $F$ maps each connected component of
%	$$
%		\leaf{}\cap\Sigma_1
%	$$
%	onto a connected component of
%	$$
%		\leaf{}\cap\Sigma_2.
%	$$
%	Hence $F$ identifies the transverse orbit foliations on $\Sigma_1$ and $\Sigma_2$ as germs at $p$.
	
Since $\Psi_*Y$ is tangent to the ambient foliation, each of its
integral curves remains in a single ambient leaf. Hence its time-one
flow preserves every ambient leaf and restricts to a biholomorphism of
germs
\[
    F:(\Sigma_1,p)\longrightarrow(\Sigma_2,p).
\]

Moreover, for every ambient leaf $\leaf{}$ meeting the neighborhood,
$F$ maps each connected component of $\leaf{}\cap\Sigma_1$ onto a
connected component of $\leaf{}\cap\Sigma_2$. Therefore, $F$ identifies
the germs at $p$ of the transverse orbit foliations on $\Sigma_1$ and
$\Sigma_2$.	
\end{proof}

\end{document}
